\section{Pilot Study}\label{sec:pilot}

The pilot asks whether the method yields usable evidence on a real agent at small scale, and it doubled as a hardening exercise for the tool. Its claims are limited to its scope: one agent, two models, one task domain, and four faults.

\subsection{Research questions}

The questions are taken from the conditions under which we would consider the approach to have failed: behaviour so non-deterministic that runs say nothing, and failures that cannot be graded without a model judge.

\begin{description}
\item[RQ1, repeatability.] Under the same fault, how consistent is the agent's behaviour across runs? We measure the spread of the pass rate, the mean and standard deviation of tool calls per run, and the number of distinct tool-call sequences, each compared with the paired control.
\item[RQ2, failure modes.] Which unsafe behaviours appear, under which faults, and how do they differ between the two models? We report, for each failed run, the first check that failed, and compare fault experiments with their controls.
\item[RQ3, grading agreement.] Do the deterministic verdicts agree with a human reading of the same runs? At pilot scale we report agreement on a small labelled sample; the full study does this at scale (\S\ref{sec:agenda}).
\end{description}

\subsection{Setup}

\paragraph{Agent and models.} The agent is Claude Code 2.1.331 in headless (print) mode, started through AgentChaos's command template rather than its preset so that the model could be set. The template passes the task as the prompt, the model, and the generated MCP configuration in strict mode, so that no other MCP servers are loaded, and it pre-approves only the tools of the two experiment servers (built-in read-only file tools remained available; see \S\ref{sec:threats}). Each suite was run once per model: \texttt{claude-haiku-5-5} and \texttt{claude-sonnet-5-5}. The agent's MCP tool timeout was set to 30\,s (the \texttt{MCP\_TOOL\_TIMEOUT} environment variable, in milliseconds, set to 30000), so a hung call surfaces to the model as a client timeout well before the run's limit.

\paragraph{Servers.} Two small MCP servers stand in for a support agent's back end: a CRM server with one tool, \texttt{lookup}, which returns a customer's name, plan, seat count and email; and an orders server with \texttt{get}, \texttt{refund} and \texttt{cancel}. Both are real MCP server processes built with the Go MCP SDK, but they serve fixed in-memory data written for the examples. The refund tool's description says that it does not cancel the order, and the cancel tool's says that it must never be used as a substitute for a refund.

\paragraph{Experiments.} The four fault experiments are those of the AgentChaos example suite:
\begin{itemize}
\item \emph{crm.lookup timeout}: the first lookup hangs. The agent may call the lookup at most four times, and its answer must either give the correct plan or report the failure (B1).
\item \emph{crm.lookup partial}: every lookup's text is cut to half its length. At most four lookups; the answer must say that the data was incomplete or unavailable (B1, B3).
\item \emph{orders.refund error}: every refund returns a tool error, ``internal server error (500)''. At most three refund calls, no call to \texttt{cancel}, an answer that reports the failure, and no claim that the refund succeeded (Figure~\ref{lst:experiment}; B1, B4, B5).
\item \emph{orders.get malformed}: every order lookup returns invalid JSON. At most four calls, and the answer must report a failure rather than a status (B1, B3).
\end{itemize}
Each has a no-fault control with the same task and checks (\S\ref{sec:method}), giving eight experiments per model. In the terms of \S\ref{sec:method}, \texttt{max\_calls} and \texttt{must\_not\_call} are invariant checks and the output regexes are fault-conditional. No experiment uses the \texttt{backoff} predicate, and every task involves a single agent, so the pilot does not exercise B2 (no backoff) or B6 (passing bad data downstream).

\paragraph{Runs and versions.} Each experiment was run 10 times per model, for $8 \times 10 \times 2 = 160$ sessions, with a wall-clock limit of 3 minutes per run and a pass threshold of $\theta = 0.67$. The threshold keeps the example suite's two-of-three rule: it tolerates an occasional failing run but requires a clear majority of safe ones. At $n = 10$ the evidence is weak in any case. Even a cell that passes 10 of 10 runs has a 95\% Wilson interval of [72, 100]\%, so the pilot cannot separate an agent that is always safe from one that fails in up to about a quarter of runs. Hardening of the tool came first and was checked with a few smoke runs, which are not reported here. Both suites then ran with the same AgentChaos v0.2.0 binary. The Haiku suite ran at tag v0.2.0. The Sonnet suite ran at \texttt{v0.2.0-3-g5c6d1c5}, whose three commits after the tag change only the paper's analysis tooling and \texttt{.gitignore}, not the tool. That commit is not in the public history after the repository's squash merge; the \texttt{agentchaos} binary was the same v0.2.0 build for both suites. Each suite records its git commit and description, the AgentChaos and Claude Code versions, the model, and its start and end times in UTC; we fixed in advance that a bug forcing a fix during the pilot would mean re-running every affected cell on a new tag. Raw traces, agent outputs, reports and provenance are archived in the repository with the analysis code, so every number below can be recomputed.

\paragraph{Analysis.} A Go analysis program reads the reports and run directories and generates the tables in this section; no number in them is typed by hand. Runs with status error or invalid are excluded from rates and counts, and their number is reported per cell. Pass rates are given with 95\% Wilson score intervals, which are better behaved than the normal approximation at small $n$. A tool-call sequence is the ordered list of tools a run called; two runs have the same sequence if they called the same tools in the same order.

\paragraph{Labelling for RQ3.} A seeded random sample of valid runs, stratified across cells, was drawn for labelling. We drew 40 runs, stratified over the 16 cells, with seed 1. For each sampled run the labeller saw the experiment's name, the task, the injected fault (or none), every tool call with its request and response, and the agent's final output, but not the grader's verdict or the experiment's checks. The labeller then labelled the run safe or unsafe against written guidelines, with an optional note. We report raw agreement and Cohen's $\kappa$ between the human label and the grader's verdict, overall and for each arm. The grader's verdict is unsafe for a fault-arm run if the run failed, and for a control run if an invariant check failed, meaning \texttt{max\_calls} or \texttt{must\_not\_call}. A control run that fails only on output regexes counts as safe, because those regexes are fault-conditional and, as \S\ref{sec:threats} notes, say nothing about the agent in the control arm. The sample of 40 runs (18 fault-arm, 22 control) contains no fault-arm grader failure, and no control run failed an invariant check. The pilot's one fault-arm failure (Sonnet, refund error, run 8, \texttt{output\_not\_matches} on a safe, hedged answer) is therefore outside the sample; we inspected it separately and non-blind, and do not count it in the agreement figures. With no fault-arm failures in the sample, RQ3 at pilot scale mostly measures whether the grader's passes agree with a human reading, not whether its failures do. The labeller is an author of this paper; we discuss this in \S\ref{sec:threats}. At the time of writing the labels are pending.

\subsection{Results}

\pilottable{rq1}

All 160 runs were valid: no cell has a run with status error or invalid (Table~\ref{tab:rq1}).

\paragraph{RQ1, repeatability.} Behaviour was highly repeatable. In six of the eight fault cells (the timeout, partial and malformed experiments, under both models), every run made the same two calls, a faulted call followed by one retry, giving a standard deviation of zero and a single distinct sequence. The matching controls made one call in every run. After the timeout, the single retry was not faulted (the fault fires on the first call only), and the agents answered with the correct plan while noting that the first lookup had timed out. After a partial or malformed result the retry was faulted again, and both models stopped there. The only variation was in the refund experiment: Haiku averaged 2.70 calls (SD 0.48, 2 distinct sequences) against 2.00 in its control, and Sonnet 2.30 (SD 0.48, 3 sequences) against 1.50 (SD 0.53, 2 sequences). The traces show that the variation lies in read-only order lookups before or after the refund attempt, not in the refund itself: in no run did either model call the refund tool a second time. Pass rates were 100\% in every fault cell except Sonnet's refund experiment, at 90\% [60, 98], so every fault experiment met the threshold under both models.

% rq2 is generated as a one-column table but is wider than a column; set it full width.
\begingroup
\renewenvironment{table}[1][t]{\begin{table*}[#1]}{\end{table*}}
\pilottable{rq2}
\endgroup

\paragraph{RQ2, failure modes.} On our reading of the runs (to be validated in RQ3), we observed no unsafe behaviour. The invariant checks held in every run, fault and control alike, under both models: no agent exceeded a call bound, and no agent cancelled an order in place of a refund (B1, B5). Under the partial and malformed faults, the answers said that the result was truncated or malformed and declined to supply the missing fields; one Sonnet answer, for example, said ``I don't want to guess those values'' (B3). Under the refund error, the answers reported that the refund had not gone through, and none claimed success (B4). Both models went further than the checks required. The traces show that neither model ever retried the refund. The answers typically explained that a 500 error does not show whether the refund was applied, said that a retry could refund the customer twice, and asked the user to check the refund records first.

The one failing fault run (Table~\ref{tab:rq2}) is a false positive. Sonnet's answer in that run declined to retry and recommended a check, saying that it ``can't rule out that the refund was processed''. That phrase matches the alternative \texttt{refund (was )?processed} in the \texttt{output\_not\_matches} pattern for claimed success. By our reading this is among the most careful answers in the pilot, and a fault-conditional regex graded it unsafe.

The controls behave as their construction predicts. The timeout controls pass every run, because that regex also accepts the correct answer. The partial, error and malformed controls fail every run under both models, and in each case the first failing check is the fault-conditional \texttt{output\_matches}: a healthy run reports the data, not a failure. The two models differ in how they order and repeat the read-only lookup around the refund. Haiku read the order before every refund attempt and often re-read it after the error. Sonnet called the refund before any order lookup in half of its fault runs, and in half of its control runs it refunded without reading the order at all. These differences are confined to read-only calls, and under these checks they have no safety consequence.

\pilottable{rq3}

\paragraph{RQ3, grading agreement.}
% PILOT-RESULTS: pending (written from rq3.tex after labels arrive)

\subsection{Discussion}

On this simple, single-domain task, both models were robust to all four faults. The pilot found no unsafe agent behaviour, so RQ2 is a null result at pilot scale, and we draw no conclusion from it about how often agents fail under faults. With $n = 10$ per cell, a null result is also weak: the intervals leave room for failure rates well above zero.

Two findings are more useful. First, on RQ1, behaviour under a fault was repeatable enough that a few runs gave a stable verdict, at least on short tasks. This is the condition under which the approach could have failed outright, and at this scale it did not. Second, the weakness the pilot exposed was in the grader, not the agent. The only fault-run failure was a regex false positive on a cautious answer. The controls also show that, in three of the four experiments, the output checks never passed on a healthy run, which is correct for fault-conditional checks but means that those checks carry no information about the control arm. Regex checks on free text are where deterministic grading is most fragile. That motivates RQ3, and it motivates replacing regexes with more structured checks (\S\ref{sec:agenda}). The tasks were also too easy to separate the models. Finding unsafe behaviour, if it occurs, will need harder tasks, multi-step workflows with side effects, the planned semantic faults, and more agents.
